% benchmark: the numbers for Part I. Every benchmark value of chapters 1 to 7 sits here.
\section{The benchmark in numbers}\label{sec:benchmark}

Part~I gives each result in words and in general symbols. This chapter puts numbers on those results at the
benchmark. The graphs of the strength path in \Cref{sec:08_coexistence,sec:09_research} use these numbers.
The subsections follow the order of Part~I. Each one points back to the chapter that it checks.

\subsection{The benchmark primitives}\label{sec:benchmark-primitives}

The benchmark has \(h=\reg{prim_h}\), \(\ell=\reg{prim_ell}\), \(\pi=\reg{prim_pi}\), \(p=\reg{prim_p}\),
\(\rho=\reg{prim_rho}\), \(c_H=\reg{prim_cH}\) and \(k=\reg{prim_k}\). Money is in units of \(\ell\), so
\(\ell=1\). The class values are point masses at \(h\) and \(\ell\). The incumbent is uniform on \([0,r]\).
The three economies of Proposition~2 have \(r_0=\reg{prim_r0}\), \(r_1=\reg{prim_r1}\) and
\(r_2=\reg{prim_r2}\). The slot has \(1+u=\eu\) and \(\beta=\reg{beta_B}\) (\Cref{sec:benchmark-slot}).
\Cref{sec:B_parameters} lists every parameter vector of the guide.

\subsection{The auction payoffs}\label{sec:benchmark-auction}

This subsection puts the benchmark values into the auction algebra of \Cref{sec:03_auction}.

\paragraph{The threshold and its derivative.} With \(h=10\), \(\ell=1\), \(p=1/2\) and \(c_H=6\), the general
threshold \eqref{eq:03-tau} and its derivative \eqref{eq:03-taup} become
\begin{equation}\label{eq:benchmark-tau}
\tau(r)=\frac{48r-3}{4(20r-r^2-1)},\qquad
\tau'(r)=\frac{3(8r^2-r+2)}{2(20r-r^2-1)^2}>0\qquad(1<r<10).
\end{equation}
At \(r=1\), \(\tau=5/8>1/2\). So the threshold lies above the prior along the whole benchmark path.

\paragraph{Derivatives at the strong economy.} At \(r_1=3\), \(\DT'=4/9\), \(g_H'=-35/72\), \(g_L'=-1/24\),
and \(\partial B_r(1/2)/\partial r=-19/72\).

\subsection{Benchmark numbers (Table 1 of the paper)}\label{sec:03_auction-table}

With \(h=10\), \(\ell=1\), \(p=1/2\) and \(c_H=6\), the formulas give Table~\ref{tab:03_auction-table1}.
Each entry is an exact fraction; the decimal is the paper's.

\begin{table}[!htbp]
\centering\small
\setlength{\tabcolsep}{4pt}
\begin{tabular}{@{}llccc@{}}
\toprule
Term & Formula & \(r_0=\reg{prim_r0}\) & \(r_1=\reg{prim_r1}\) & \(r_2=\reg{prim_r2}\)\\
\midrule
\(t_0\) & \(1/2-1/(4r)\) & \(7/24=\reg{p2_base_r0_t0}\) & \(5/12=\reg{p2_base_r1_t0}\) & \(31/72=\reg{p2_base_r2_t0}\)\\
\(t_H\) & \(r/2+1/(8r)\) & \(169/240=\reg{p2_base_r0_th}\) & \(37/24=\reg{p2_base_r1_th}\) & \(1321/720=\reg{p2_base_r2_th}\)\\
\(t_L\) & \(1-3/(8r)\) & \(11/16=\reg{p2_base_r0_tl}\) & \(7/8=\reg{p2_base_r1_tl}\) & \(43/48=\reg{p2_base_r2_tl}\)\\
\(\DT\) & \((r-1)^2/(2r)\) & \(1/60=\reg{Delta_r0}\) & \(2/3=\reg{Delta_r1}\) & \(169/180=\reg{Delta_r2}\)\\
\(g_H\) & \(10-t_H\) & \(2231/240=\reg{p2_base_r0_gh}\) & \(203/24=\reg{p2_base_r1_gh}\) & \(5879/720=\reg{p2_base_r2_gh}\)\\
\(g_L\) & \(3/(8r)\) & \(5/16=\reg{p2_base_r0_gl}\) & \(1/8=\reg{p2_base_r1_gl}\) & \(5/48=\reg{p2_base_r2_gl}\)\\
\(B_r(1/2)\) & \(5-r/4+1/(8r)\) & \(1153/240=\reg{Bprior_r0}\) & \(103/24=\reg{Bprior_r1_B}\) & \(2977/720=\reg{Bprior_r2}\)\\
\(\tau(r)\) & \eqref{eq:benchmark-tau} & \(195/308=\reg{tau_r0}\) & \(141/200=\reg{tau_r1}\) & \(4245/5804=\reg{tau_r2}\)\\
\bottomrule
\end{tabular}
\caption{Acquisition payoffs in the benchmark \paper{Table~1}. They do not depend on the slot.}
\label{tab:03_auction-table1}
\end{table}

The row for \(B_r(1/2)\) uses \((g_H+g_L)/2=\tfrac12\bigl(10-\tfrac r2-\tfrac1{8r}+\tfrac3{8r}\bigr)
=5-\tfrac r4+\tfrac1{8r}\).

Note the size of the change from \(r_0\) to \(r_1\). The spread rises \(40\) times
(\((2/3)/(1/60)=40\)). Gross profit at the prior falls by \(41/80\), about \(10.7\) percent.
% verify/03_auction.py
The threshold rises from \reg{tau_r0} to \reg{tau_r1}, as \eqref{eq:03-taup} requires.

At \(r_1=3\), entry adds \(w_L=t_L-t_0=11/24\) in class \(L\) and \(w_H=t_H-t_0=9/8\) in class \(H\).
The surplus split holds: \(\DT+g_H-g_L=2/3+203/24-1/8=16/24+203/24-3/24=216/24=9=h-\ell\).

\paragraph{Proposition 1 off the benchmark.} The script \texttt{verify/03\_auction.py} checks \paper{(8)} for a
non-uniform \(F(y)=(y/\bar r)^2\) and classes with an atom: the formula and a direct double expectation
agree to twelve digits.

\paragraph{The exclusion-reserve example.} Set \(p=\reg{rsv_points_excl}\) in the example of
\Cref{sec:03_auction-prop1c}. The spread falls for \(r<\sqrt3\,p\approx1.749\). % verify/03_auction.py
So from \(r=6/5\) to \(r=3/2\) the spread \emph{falls}, from \(6921/8000=0.865125\) to
\(7601/10000=0.7601\). % verify/03_auction.py
It rises again later and reaches \(20001/20000\) at \(r=3\). % verify/03_auction.py

\subsection{The slot and the beliefs}\label{sec:benchmark-slot}

This section gives the numbers for the slot algebra of \Cref{sec:04_slot_beliefs}.

\paragraph{An order law with wrong-signed orders.}
Take $\lambda=1/3$, $\beta=1/4$, and $\pi=2/5$.
Let the profiles on $(-1,0,1)$ be
\[
\sigma_H=(1/5,3/10,1/2),\qquad \sigma_L=(3/5,3/10,1/10).
\]
Both profiles include a wrong-signed order.
Direct substitution gives
\[
f_H=(7/30,13/30,1/3),\qquad f_L=(11/30,13/30,1/5).
\]
The posteriors are $(14/47,2/5,10/19)$.
The same formulas apply before we ask whether these orders are optimal.

\paragraph{The zero-order bound fails when \texorpdfstring{$\beta>1/3$}{beta > 1/3}.}
Let class $H$ always abstain and class $L$ always sell.
Take the exact example $\lambda=6/17$, $\beta=2/5$, and $\pi=1/2$.
The masses and ratios are
\[
b=\frac{22}{85},\quad b_0=\frac{11}{85},\quad
\lambda=\frac{30}{85},\quad
1+u=\frac{26}{11},\quad \frac{f_H(0)}{f_L(0)}=\frac{41}{11}.
\]
Consequently,
\[
M=\frac{26}{37},\qquad \mu(0)=\frac{41}{52},\qquad
\mu(0)-M=\frac{165}{1924}>0.
\]

\paragraph{The benchmark slot.}
In the benchmark, $\beta=\reg{beta_B}$ and $1+u=\eu$.
With this choice, the top posterior equals the Laplace value at scale $2$ (\Cref{sec:10_robustness-sizes}).
The presence probability is derived, rather than separately rounded:
\[
\lambda=\frac{u\beta}{1+u\beta}=\reg{lambda_B}.
\]
Thus $b=\reg{b_B}$ and $b+\lambda=\reg{b_plus_lambda_B}$.
At $\pi=\reg{prim_pi}$,
\[
m=\frac1{1+\eu}=\reg{m_half_B},\qquad
M=\frac{\eu}{1+\eu}=\reg{M_half_B}.
\]
The script keeps the full value of $\lambda$ throughout the calculation.

\paragraph{Prices at the strong benchmark.}
At $r_1=\reg{prim_r1}$, $t_0=5/12$, $t_L-t_0=11/24$, and $\DT=2/3$.
The threshold is $\tau=141/200=\reg{tau_r1}$.
Under full orders, the benchmark gives
\[
\begin{array}{c|ccc}
x&-1&0&1\\ \hline
\mu(x)&\reg{m_half_B}&1/2&\reg{M_half_B}\\
e(x)&1/4&1/4&1\\
P(x)&\reg{price_r1_m1_B}&\reg{price_r1_0_B}&\reg{price_r1_p1_B}
\end{array}
\]
Only the purchase posterior exceeds $\tau$.
The script substitutes these prices into the inverse formula and recovers all three posteriors.

Run \texttt{python3 slot/study\_guide/verify/04\_slot\_beliefs.py} from the repository root.
It prints the order laws, the zero-order counterexample, the posterior bounds, and the price inversion.

\subsection{The investor's orders at the benchmark}\label{sec:benchmark-investor}

These values check the trading bounds of \Cref{sec:05_investor}.

\paragraph{A mixed candidate inside the bounds.}
For a concrete candidate that is not an equilibrium, use the mixed profile of \Cref{sec:benchmark-slot} at $r=3$, where $\DT=2/3$.
Every posterior is below $\tau=141/200$, so set $e\equiv\rho=1/4$.
The three posteriors differ, so the three prices differ and the price beliefs equal the order posteriors.
Thus (E1)--(E4) and (E6) hold, although the mixed orders are not optimal.
Here $u=2$, $m(2/5)=2/11$, and $1-M(2/5)=1/3$, so $\kappa=2/11$.
The exact gross returns are
\[
\rho\kappa\DT=\frac14\cdot\frac2{11}\cdot\frac23=\frac1{33},\qquad
K_H=\frac14\cdot\frac23\left(1-\frac{10}{19}\right)=\frac3{38},\qquad
K_L=\frac14\cdot\frac23\cdot\frac{14}{47}=\frac7{141}.
\]
Each lies between $1/33$ and $2/3$, as the bound requires without investor optimality.

\paragraph{The floor factor at other priors.}
At $1+u=\eu$, $\kappa(1/2)=\reg{m_half_B}$.
At $\pi=1/4$ and $3/4$, it equals approximately $0.109232$. % verify/05_investor.py

\paragraph{The trading-cost window at the benchmark.}
At the weak benchmark, $\DT(r_0)=1/60$ and $k=1/50$.
Hence the largest possible net return is at most
\[
\frac1{60}-\frac1{50}=-\frac1{300}<0.
\]
At the strong benchmark,
\[
\rho\kappa\DT(r_1)-k
=\frac1{6(1+\eu)}-\frac1{50}
=\reg{margin_kstrong_B}>0.
\]
In the full-order candidate, $e(1)=1$ and $e(-1)=1/4$.
Thus $K_H=2/[3(1+\eu)]$ and $K_L=1/[6(1+\eu)]$.
Their net returns are $\reg{p2_base_r1_choice_h}$ and $\reg{p2_base_r1_choice_l}$.

Run \texttt{python3 slot/study\_guide/verify/05\_investor.py} from the repository root.
It prints the six payoffs in each benchmark economy, the return bounds, and the prior symmetry.
It also prints the mixed candidate's exact returns.

\subsection{The main result at the benchmark}\label{sec:benchmark-main}

This subsection evaluates the steps of \Cref{sec:06_main_result} at the benchmark.

\subsubsection{Entry and ownership}\label{sec:benchmark-entry}
At $r_1=\reg{prim_r1}$, \Cref{tab:03_auction-table1} gives $g_L=1/8$ and $g_H=203/24$, so
$\tau_1=(6-1/8)/(25/3)=141/200$.
Since $\pi<\tau_1<M$, only the purchase recruits.
With $\beta=\reg{beta_B}$ and $u=\eu-1$, the script gives
\begin{align*}
\lambda&=\reg{lambda_B},& b&=\reg{b_B},& b+\lambda&=\reg{b_plus_lambda_B},\\
e_H&=\reg{eH_r1_B},& e_L&=\reg{eL_r1_B},& \Ent_1&=\reg{entry_r1_B}.
\end{align*}
Finally $\Own_{H,1}=\pi e_H=\reg{own_r1_B}$.
The weak values are $\Ent_0=\reg{entry_r0}$ and $\Own_{H,0}=\reg{own_r0}$.

\subsubsection{The very strong economy and the entry ceiling}\label{sec:benchmark-ceiling}
At $r_2=\reg{prim_r2}$, gross profit at the top posterior is
\[
B_{r_2}(M)=\frac5{48}+\frac{1451}{180}M
=\reg{p2_base_r2_low_profit_plus}<6.
\]
The cost margin is $\reg{margin_return_B}$.
Entry is $\rho$, yet the three prices still differ (\Cref{tab:07-outcomes}).

For the uniform benchmark, expand gross profit at the top posterior:
\begin{align*}
B_r(M)
&=(1-M)\frac{\ell^2-p^2}{2r}
  +M\left(h-\frac r2-\frac{p^2}{2r}\right)\\
&=Mh-\frac{Mr}{2}+\frac{(1-M)\ell^2-p^2}{2r}.
\end{align*}
Each class profit falls strictly with $r$, so this weighted average also falls strictly.
The entry ceiling $r_C$ solves $B_{r_C}(M)=c_H$.
Multiplication by $2r_C$ gives
\[
Mr_C^2-2(Mh-c_H)r_C-[(1-M)\ell^2-p^2]=0.
\]
In the benchmark, $(1-M)\ell^2-p^2=1/(1+\eu)-1/4>0$, so the product of the two roots, $-[(1-M)\ell^2-p^2]/M$, is negative.
The positive root is
\[
r_C=\frac{Mh-c_H+\sqrt{(Mh-c_H)^2+M[(1-M)\ell^2-p^2]}}{M}
=\reg{rC_B}.
\]

\subsubsection{The admissible band at the prior one half}\label{sec:benchmark-band}
Fix the benchmark auction, costs, and strengths, and let $\pi=1/2$.
The threshold $\tau(r)=(48r-3)/[4(20r-r^2-1)]$ of \Cref{sec:03_auction-derivatives} gives
\[
\tau(r_0)=\frac{195}{308},\qquad
\tau(r_1)=\frac{141}{200},\qquad
\tau(r_2)=\frac{4245}{5804},
\]
and \Cref{tab:03_auction-table1} gives $\Delta_0=1/60$, $\Delta_1=2/3$, and $\Delta_2=169/180$.
The weak-economy margins do not depend on the slot:
\[
c_H-B_0(1/2)=6-\frac{1153}{240}=\frac{287}{240}>0,\qquad
k-\Delta_0=\frac1{50}-\frac1{60}=\frac1{300}>0.
\]
The right side of (H1) becomes
\begin{align*}
M>\frac{141}{200}
&\iff 200(1+u)>141(2+u)\\
&\iff 59u>82
\iff u>\frac{82}{59}.
\end{align*}
The ceiling (H3) becomes
\begin{align*}
M<\frac{4245}{5804}
&\iff5804(1+u)<4245(2+u)\\
&\iff1559u<2686
\iff u<\frac{2686}{1559}.
\end{align*}
The right side of (H2) becomes
\begin{align*}
\frac1{50}<\frac1{6(2+u)}
&\iff6(2+u)<50
\iff u<\frac{19}{3}.
\end{align*}
Equivalently, it requires $M<22/25$.
Since $2686/1559<19/3$, the trading bound is redundant within the narrower band.
We therefore obtain the exact equivalence
\[
\text{(H1)--(H3)}
\iff \frac{82}{59}<u<\frac{2686}{1559}
\iff \frac{141}{200}<M<\frac{4245}{5804}.
\]
The equivalent presence interval is
\[
\frac{82\beta}{59+82\beta}<\lambda<\frac{2686\beta}{1559+2686\beta}.
\]
At $\beta=\reg{beta_B}$, its endpoints are $\reg{lambda_lo_B}$ and $\reg{lambda_hi_B}$.
The benchmark $\lambda=\reg{lambda_B}$ lies strictly between them.

\subsubsection{The prior window and the boundary cases}\label{sec:benchmark-priors}
Now fix $1+u=\eu$ and vary the prior.
The right side of (H1) and (H3) give
\[
\frac{141}{141+59\eu}<\pi<\frac{4245}{4245+1559\eu}.
\]
The endpoints are approximately $\reg{prior_lo_B}$ and $\reg{prior_hi_B}$.
The left side of (H1) requires $\pi<\tau(r_0)=195/308$.
The right side of (H2) requires
\[
\kappa(\pi)>\frac{k}{\rho\Delta_1}=\frac{1/50}{(1/4)(2/3)}=\frac3{25}.
\]
Both terms in the minimum must exceed $3/25$:
\begin{align*}
m(\pi)>\frac3{25}
&\iff 25\pi>3[\pi+(1-\pi)\eu]\\
&\iff\pi>\frac{3\eu}{22+3\eu},\\
1-M(\pi)>\frac3{25}
&\iff25(1-\pi)>3[1-\pi+\pi\eu]\\
&\iff\pi<\frac{22}{22+3\eu}.
\end{align*}
This interval has endpoints $\reg{prior_kappa_lo_B}$ and $\reg{prior_kappa_hi_B}$.
It contains both the full prior window and the window for parts (i) and (ii).
The full intersection therefore equals the first prior interval above.
Without (H3), parts (i) and (ii) require exactly
\[
\frac{141}{141+59\eu}<\pi<\frac{195}{308}.
\]

The boundary cases at the prior one half read as follows.
For $u<82/59$, (H2) still forces full orders at $r_1$, but $M<\tau(r_1)$.
For $2686/1559\le u<19/3$, full orders persist at both stronger strengths, and a purchase at $r_2$ recruits under C1.
For $u\ge19/3$, the right side of (H2) fails.

At the lower boundary $u=82/59$, the posteriors satisfy $M=141/200$ and $m=59/200$.
Orders at $r_0$ remain uniquely zero because $\Delta_0-k=-1/300$.
Orders at $r_1$ and $r_2$ remain uniquely full because
\[
\rho\kappa\Delta_1-k
=\frac14\frac{59}{200}\frac23-\frac1{50}=\frac7{240}>0.
\]
At $\beta=\reg{beta_B}$, the boundary value of $\Ent_1$ under C1 is approximately $0.515737$. % verify/06_main_result.py
It differs from the benchmark value because the boundary slot has a different $u$ and $\lambda$.

\subsubsection{Why H1 needs the weak prior, and the benchmark margins}\label{sec:benchmark-scope}
Keep the benchmark slot and other primitives, but set $c_H=\reg{nec_cH}$ and $r_2=\reg{nec_r2}$.
The proposed weaker list has positive margins:
\begin{align*}
c_H-B_1(1/2)&=\frac{49}{120}>0,& B_1(M)-c_H&=\reg{nec_margin_recruit_B}>0,\\
k-\Delta_0&=\frac1{300}>0,& \rho\kappa\Delta_1-k&=\reg{nec_margin_kstrong_B}>0,\\
c_H-B_2(M)&=\reg{nec_margin_return_B}>0.
\end{align*}
But the omitted weak-prior comparison has the opposite sign:
\[
c_H-B_0(1/2)=\frac{47}{10}-\frac{1153}{240}=-\frac5{48}<0.
\]
Entry at $r_0$ is one, and entry at $r_1$ is $\reg{entry_r1_B}$.
Entry falls, so the intended comparison fails.

At the benchmark itself, the five positive margins of (H1)--(H3) are
\[
\begin{array}{c|c}
c_H-B_0(\pi)&\reg{margin_weakcost_B}\\
B_1(M)-c_H&\reg{margin_recruit_B}\\
k-\Delta_0&\reg{margin_kweak_B}\\
\rho\kappa\Delta_1-k&\reg{margin_kstrong_B}\\
c_H-B_2(M)&\reg{margin_return_B}
\end{array}
\]
These margins prove that the set is nonempty and contains a neighborhood of the benchmark.
The small margin at the ceiling measures sensitivity, rather than an equality.

Run \texttt{python3 slot/study\_guide/verify/06\_main\_result.py} from the repository root.
It prints the three tuples, auction quantities, entry, ownership, exact band, prior window, ceiling, and H1 counterexample.

\subsection{The controls at the benchmark}\label{sec:benchmark-controls}

This section puts the benchmark numbers into the theorem and the two controls of \Cref{sec:07_controls}.

\subsubsection{The two examples in numbers}

\texttt{verify/07\_controls.py} recomputes both examples with exact fractions.

\paragraph{Example A (breaks P2).} Take \(p=1/2\), point-mass
classes \(12\) and \(11\) with equal priors, \(\rho=1/4\) and \(c_H=6\). Let
\(F_0=\tfrac12U[0,1]+\tfrac12U[8,9]\) and \(F_1=\tfrac12U[0,1]+\tfrac12U[9,10]\). The upper component moves up, so
\(F_1\le F_0\). In both economies \(\mu(S)=1/2\) and \(\E[\theta\mid S]=23/2\). The expected payment is
\(\tfrac12\cdot\tfrac12+\int_{1/2}^1x\,dx=\tfrac58\) on \(U[0,1]\), \(\tfrac{17}2\) on \(U[8,9]\) and \(\tfrac{19}2\) on
\(U[9,10]\).
\begin{center}\small
\begin{tabularx}{\linewidth}{@{}lXX@{}}
\toprule
 & Economy 0: signal reports the component & Economy 1: signal independent of \(R\)\\
\midrule
Gross profit after the signal & \(\tfrac{23}2-\tfrac58=\tfrac{87}8\) or \(\tfrac{23}2-\tfrac{17}2=3\) &
\(\tfrac{23}2-\tfrac12\bigl(\tfrac58+\tfrac{19}2\bigr)=\tfrac{103}{16}\) after either signal\\
High-cost preparation & after the low component only & after every signal\\
Entry & \(\tfrac14+\tfrac34\cdot\tfrac12=\tfrac58\) & \(1\)\\
\bottomrule
\end{tabularx}
\end{center}
Unconditional gross profit falls from \(111/16\) to \(103/16\). In economy~0 the profit after the strong component,
\(3\), fails \(c_H\). In economy~1 the average profit \(103/16\) clears \(c_H=6\).

\paragraph{Example B (breaks P1).} Take
\(G_L=\tfrac12\delta_1+\tfrac12\delta_3\), \(G_H=\tfrac12\delta_2+\tfrac12\delta_4\), \(p=1/2\), \(\rho=1/4\) and
\(c_H=21/10\). Economy~0 has \(R\sim U[0,1/2]\). Economy~1 has \(R\sim U[0,3/5]\). The expected payment is
\(\tfrac12\) on \(U[0,\tfrac12]\), and
\(\tfrac12\cdot\tfrac56+\int_{1/2}^{3/5}\tfrac{x}{3/5}\,dx=\tfrac5{12}+\tfrac{11}{120}=\tfrac{61}{120}\) on \(U[0,\tfrac35]\).
\begin{center}\small
\begin{tabularx}{\linewidth}{@{}lXX@{}}
\toprule
 & Economy 0: \(\theta\) independent of \(S\) & Economy 1: \(S\) selects the value\\
\midrule
\(\E[\theta\mid S]\) & \(\tfrac12\cdot2+\tfrac12\cdot3=\tfrac52\) & \(\tfrac32\) (lower) or \(\tfrac72\) (upper)\\
Gross profit after the signal & \(\tfrac52-\tfrac12=2<\tfrac{21}{10}\) &
\(\tfrac32-\tfrac{61}{120}=\tfrac{119}{120}\) or \(\tfrac72-\tfrac{61}{120}=\tfrac{359}{120}\)\\
Entry & \(\tfrac14\) & \(\tfrac14+\tfrac34\cdot\tfrac12=\tfrac58\)\\
\bottomrule
\end{tabularx}
\end{center}
The class gross profits move from \((g_L,g_H)=(3/2,5/2)\) to \((179/120,299/120)\). After the upper signal, gross
profit \(359/120\) exceeds \(c_H=21/10\).

\subsubsection{Control 1 in numbers}

At the benchmark slot, \(m=\reg{m_half_B}\) and \(M=\reg{M_half_B}\). The threshold starts at \(\tau(1)=5/8>1/2\), and
the entry ceiling is \(r_C\approx\reg{rC_B}\).

A purchase has probability \(b+\lambda/2=\reg{d_slot_B}\). So frozen entry is
\(\tfrac14+\tfrac34\times\reg{d_slot_B}=\reg{p4_frozen_entry_below_ceiling}\) for \(1<r\le r_C\), and
\(\rho=\reg{p4_frozen_entry_above_ceiling}\) for \(r>r_C\).

The crossing probabilities:
\begin{itemize}[leftmargin=2em]
\item \(r_0\to r_1\): \(B_{r_1}(M)=\reg{p2_base_r1_low_profit_plus}\ge c_H\) and \(B_{r_0}(1/2)=\reg{Bprior_r0}<c_H\).
The probability is \(\reg{p4_control_crossing_r0_r1}\).
\item \(r_1\to r_2\): \(B_{r_2}(M)=\reg{p2_base_r2_low_profit_plus}<c_H\). The probability is
\(\reg{p4_control_crossing_r1_r2}\), exactly the fall in entry.
\end{itemize}

At \(r_0\), \(B_{r_0}(M)-c_H=\reg{margin_p4_weak_top_B}>0\). The informed orders earn
\begin{align*}
\text{purchase: }&\DT(r_0)(1-M)-k=\tfrac1{60}\times\reg{m_half_B}-\tfrac1{50}=\reg{p4_frozen_r0_buy_tp_payoff},\\
\text{sale: }&\rho\,\DT(r_0)\,m-k=\tfrac14\times\tfrac1{60}\times\reg{m_half_B}-\tfrac1{50}=\reg{p4_frozen_r0_sell_tm_payoff}.
\end{align*}
Thus frozen entry stays at \reg{p4_frozen_r0_entry} through \(r_C\), then falls by \reg{p4_control_crossing_r1_r2}.

\subsubsection{Control 2 in numbers}

At \(r_0\), \(\DT(r_0)=1/60<k=1/50\). At \(r_1\),
\(\rho m\DT(r_1)=\tfrac14\times\reg{m_half_B}\times\tfrac23=0.044824>k\).
% verify/07_controls.py
The sale margin is \(\reg{margin_p4_strong_trade_B}\), as in the equilibrium at \(r_1\).
At \(r_1\) the prices are \(\reg{p4_hidden_r1_sell_price}\), \(\reg{p4_hidden_r1_zero_price}\) and
\(\reg{p4_hidden_r1_buy_price}\). Entry is
\(\reg{p4_hidden_r0_entry}\), ownership \(\rho/2=\reg{p4_hidden_r1_own}\), and proceeds \(t_0+\rho\,w(1/2)\): at \(r_1\),
\(\tfrac5{12}+\tfrac14\bigl(\tfrac78-\tfrac5{12}+\tfrac13\bigr)=\tfrac{59}{96}=\reg{p4_hidden_r1_proceeds}\).

\subsubsection{The outcomes table decoded}\label{sec:07_controls-table}

Table~\ref{tab:07-outcomes} collects the comparisons of Table~3 in the paper.
\Cref{sec:03_auction-table} gives the auction inputs.
\Cref{sec:06_main_result-entry} derives equilibrium entry and ownership.
The controls use the same aggregation formulas.

\begin{table}[ht]
\centering\small
\caption{Equilibrium outcomes and information controls \paper{Table~3}.}\label{tab:07-outcomes}
\begin{tabular}{@{}llccc@{}}
\toprule
Regime & Strength & Entry & High-class ownership & Target proceeds\\
\midrule
Equilibrium & \(r_0\) & \reg{entry_r0} & \reg{own_r0} & \reg{p4_base_r0_proceeds}\\
& \(r_1\) & \reg{entry_r1_B} & \reg{own_r1_B} & \reg{p4_base_r1_proceeds}\\
& \(r_2\) & \reg{entry_r2} & \reg{own_r2_B} & \reg{p4_base_r2_proceeds}\\
Frozen full orders & \(r_0\) & \reg{p4_frozen_r0_entry} & \reg{p4_frozen_r0_own} & \reg{p4_frozen_r0_proceeds}\\
& \(r_1\) & \reg{p4_frozen_r1_entry} & \reg{p4_frozen_r1_own} & \reg{p4_frozen_r1_proceeds}\\
Hidden price & \(r_0\) & \reg{p4_hidden_r0_entry} & \reg{p4_hidden_r0_own} & \reg{p4_hidden_r0_proceeds}\\
& \(r_1\) & \reg{p4_hidden_r1_entry} & \reg{p4_hidden_r1_own} & \reg{p4_hidden_r1_proceeds}\\
\bottomrule
\end{tabular}
\end{table}

\paragraph{Expected proceeds.} \(R_T=t_0+\tfrac12e_Hw_H+\tfrac12e_Lw_L\). At \(r_1\):
\(\tfrac5{12}+\tfrac12\times\reg{eH_r1_B}\times\tfrac98+\tfrac12\times\reg{eL_r1_B}\times\tfrac{11}{24}=\reg{p4_base_r1_proceeds}\).
At \(r_0\) and \(r_2\), \(e_H=e_L=\rho\), so \(R_T=t_0+\rho w(1/2)\).
The values are \(377/960\) and \(1913/2880\). As a check,
competitive prices make \(R_T=\E[P(X)]\). With \(e\equiv\rho\), \(w\) is linear and the posteriors average to \(1/2\), so
\(R_T\) equals the zero-order price.

\paragraph{The control rows.} The frozen tuple at \(r_0\) has the same \(e_H\) and \(e_L\) as at \(r_1\), because the same
purchase recruits; so entry and ownership repeat the \(r_1\) values. Proceeds use the \(r_0\) inputs:
\(\tfrac7{24}+\tfrac12\times\reg{eH_r1_B}\times\tfrac{33}{80}+\tfrac12\times\reg{eL_r1_B}\times\tfrac{19}{48}=\reg{p4_frozen_r0_proceeds}\).
At \(r_1\) the frozen tuple is the equilibrium. At \(r_2\), frozen full orders give the equilibrium row, since \(r_2>r_C\).
The hidden-price control at \(r_0\) also gives the equilibrium row, since its price carries no information.

\subsection{The main result in numbers}\label{sec:benchmark-headline}

Table~\ref{tab:benchmark-headline} puts the numbers on Table~\ref{tab:01_brief-numbers}. It compares the three
economies of Proposition~2 at the benchmark.

\begin{table}[htbp]
\centering\small
\begin{tabular}{@{}lccc@{}}
\toprule
 & Weak (\(r_0=\reg{prim_r0}\)) & Strong (\(r_1=\reg{prim_r1}\)) & Very strong (\(r_2=\reg{prim_r2}\))\\
\midrule
Gross profit at the prior, \(B_r(1/2)\) & \reg{Bprior_r0} & \reg{Bprior_r1_B} (lower) & \reg{Bprior_r2}\\
Target-payoff spread, \(\DT\) & \reg{Delta_r0} & \reg{Delta_r1} (\(40\times\)) & \reg{Delta_r2}\\
Preparation threshold, \(\tau(r)\) & \reg{tau_r0} & \reg{tau_r1} & \reg{tau_r2}\\
Top posterior, \(M(1/2)\) & \reg{M_half_B} & \reg{M_half_B} & \reg{M_half_B}\\
Investor orders, classes \(H\) and \(L\) & \((0,0)\) & \((1,-1)\) & \((1,-1)\)\\
Price informative? & no & yes & yes\\
High-cost preparation & none & after a purchase & none\\
Entry \(\Ent\) & \reg{entry_r0} & \textbf{\reg{entry_r1_B}} & \reg{entry_r2}\\
High-class ownership \(\Own_H\) & \reg{own_r0} & \textbf{\reg{own_r1_B}} & \reg{own_r2_B}\\
\bottomrule
\end{tabular}
\caption{The main result at the benchmark \paper{Tables~1 and~3}.}
\label{tab:benchmark-headline}
\end{table}

From \(r_0\) to \(r_1\), gross profit falls by about \(10.7\) percent, while the spread rises \(40\) times. % verify/03_auction.py Entry more than doubles, by a factor of
\(2.09\). % verify/01_brief.py
At \(r_2\) the price stays informative, but \(\tau(r_2)=\reg{tau_r2}\) lies just above
\(M(1/2)=\reg{M_half_B}\). So even a purchase cannot recruit a high-cost challenger, and entry falls back
to \(\rho\). Entry is \emph{non-monotone} in incumbent strength.

\paragraph{Why entry doubles while profit falls.}
Profit falls at \emph{every} belief, but the belief that matters changes. At \(r_0\) the price says nothing,
so the challenger decides at the prior \(1/2\), and \(B_{r_0}(1/2)=\reg{Bprior_r0}<c_H=6\). At \(r_1\) a
purchase moves the belief to \(\reg{M_half_B}\), and \(B_{r_1}(M)\) exceeds \(6\) by
\(\reg{margin_recruit_B}\). A purchase occurs with probability \(b+\lambda/2\) at the prior \(1/2\), so entry is
\(\rho+(1-\rho)(b+\lambda/2)=\reg{entry_r1_B}\).
